% Lecture 12 - Page 1 (3-Jan-1995).
% Fragment for % Lecture 12 - Page 1 (3-Jan-1995).
% Fragment for % Lecture 12 - Page 1 (3-Jan-1995).
% Fragment for % Lecture 12 - Page 1 (3-Jan-1995).
% Fragment for \input{courses/pages/curs12-pagina1-en}.
% Faithful translation retaining the manuscript's sequence and notation.
% No preamble; start a new lecture on a new page.

\clearpage
\subsection*{Lecture 12 - Page 1}

\begin{center}
\fbox{\textbf{LECTURE 12}}
\end{center}
\begin{flushright}
\textit{3-Jan-1995}
\end{flushright}

\medskip
\noindent\underline{Generalities on monoids}

\bigskip
\noindent\fbox{T.}\quad Let $(M,\cdot,e)$ be a monoid, $\varnothing\neq S\subseteq M$,
\[
\langle S\rangle
=\{x_1\cdots x_n\mid x_i\in S,\ n\geq 1\}\cup\{e\}.
\]
\noindent
\begin{enumerate}
\item[(1)] $\langle S\rangle$ is a submonoid of $M$, $\langle S\rangle\leq M$.
\item[(2)] $N\leq M$, $S\subseteq N\ \Longrightarrow\ \langle S\rangle\subseteq N$.
\end{enumerate}

\noindent A non-empty subset $S\subseteq M$ is called a \underline{submonoid}
if $S$ is closed under the operation and $e\in S$.

\bigskip
\noindent\underline{Proof.} Evident.

\bigskip
\noindent\underline{Remark.}\quad
$\langle S\rangle$ is the smallest submonoid of $M$ containing $S$.

\medskip
\noindent $\langle S\rangle$ is called the submonoid of $M$ generated by $S$.

\bigskip
\noindent\underline{Definition.}\quad
Let $M$ be a monoid. A subset $S\subseteq M$ is called a
\underline{generating set} for $M$ if
\[
\langle S\rangle=M.
\]

\bigskip
\noindent\underline{Example.}\quad
For $(\mathbb{N}^{*},\cdot,1)$, a generating set is the set of prime numbers,
or the set of prime numbers together with $\{1\}$.
% The handwritten example reads “{prime numbers} or {prime numbers}+{1}”.
% The final marks at the line's end are unclear; no additional example supplied.

\bigskip
\noindent\underline{Definition.}\quad
A monoid $M$ is called \underline{finitely generated}
if its generating set is finite.

\medskip
\noindent\underline{Remark.}\quad $M=\langle M\rangle$.

\bigskip
\noindent $R$ - a congruence on $M$
(an equivalence relation compatible with $\cdot$).
\[
\frac{M}{R}=\{\hat a\mid a\in M\},
\qquad \hat a\,\hat b=\widehat{ab},
\qquad \left(\frac{M}{R},\cdot,\hat e\right)
\text{ - the quotient monoid.}
\]

\bigskip
\noindent\underline{Example.}\quad
$(\mathbb{Z},\cdot,1)$, $R_n$ - congruence modulo $n$,
\[
\frac{\mathbb{Z}}{R_n}=(\mathbb{Z}_n,\cdot,\hat 1).
\]
% The manuscript uses equality here to identify the quotient with the quotient monoid.

\bigskip
\noindent\underline{Definition.}\quad
Let $(M,\cdot,e)$ and $(M',\cdot,e')$ be two monoids.
A map $f:M\to M'$ is called a \underline{monoid homomorphism} if:
\begin{enumerate}
\item[(1)] $f(xy)=f(x)f(y)$, $(\forall)\,x,y\in M$.
\item[(2)] $f(e)=e'$.
\end{enumerate}
% The definition ends at the bottom of the manuscript page.
.
% Faithful translation retaining the manuscript's sequence and notation.
% No preamble; start a new lecture on a new page.

\clearpage
\subsection*{Lecture 12 - Page 1}

\begin{center}
\fbox{\textbf{LECTURE 12}}
\end{center}
\begin{flushright}
\textit{3-Jan-1995}
\end{flushright}

\medskip
\noindent\underline{Generalities on monoids}

\bigskip
\noindent\fbox{T.}\quad Let $(M,\cdot,e)$ be a monoid, $\varnothing\neq S\subseteq M$,
\[
\langle S\rangle
=\{x_1\cdots x_n\mid x_i\in S,\ n\geq 1\}\cup\{e\}.
\]
\noindent
\begin{enumerate}
\item[(1)] $\langle S\rangle$ is a submonoid of $M$, $\langle S\rangle\leq M$.
\item[(2)] $N\leq M$, $S\subseteq N\ \Longrightarrow\ \langle S\rangle\subseteq N$.
\end{enumerate}

\noindent A non-empty subset $S\subseteq M$ is called a \underline{submonoid}
if $S$ is closed under the operation and $e\in S$.

\bigskip
\noindent\underline{Proof.} Evident.

\bigskip
\noindent\underline{Remark.}\quad
$\langle S\rangle$ is the smallest submonoid of $M$ containing $S$.

\medskip
\noindent $\langle S\rangle$ is called the submonoid of $M$ generated by $S$.

\bigskip
\noindent\underline{Definition.}\quad
Let $M$ be a monoid. A subset $S\subseteq M$ is called a
\underline{generating set} for $M$ if
\[
\langle S\rangle=M.
\]

\bigskip
\noindent\underline{Example.}\quad
For $(\mathbb{N}^{*},\cdot,1)$, a generating set is the set of prime numbers,
or the set of prime numbers together with $\{1\}$.
% The handwritten example reads “{prime numbers} or {prime numbers}+{1}”.
% The final marks at the line's end are unclear; no additional example supplied.

\bigskip
\noindent\underline{Definition.}\quad
A monoid $M$ is called \underline{finitely generated}
if its generating set is finite.

\medskip
\noindent\underline{Remark.}\quad $M=\langle M\rangle$.

\bigskip
\noindent $R$ - a congruence on $M$
(an equivalence relation compatible with $\cdot$).
\[
\frac{M}{R}=\{\hat a\mid a\in M\},
\qquad \hat a\,\hat b=\widehat{ab},
\qquad \left(\frac{M}{R},\cdot,\hat e\right)
\text{ - the quotient monoid.}
\]

\bigskip
\noindent\underline{Example.}\quad
$(\mathbb{Z},\cdot,1)$, $R_n$ - congruence modulo $n$,
\[
\frac{\mathbb{Z}}{R_n}=(\mathbb{Z}_n,\cdot,\hat 1).
\]
% The manuscript uses equality here to identify the quotient with the quotient monoid.

\bigskip
\noindent\underline{Definition.}\quad
Let $(M,\cdot,e)$ and $(M',\cdot,e')$ be two monoids.
A map $f:M\to M'$ is called a \underline{monoid homomorphism} if:
\begin{enumerate}
\item[(1)] $f(xy)=f(x)f(y)$, $(\forall)\,x,y\in M$.
\item[(2)] $f(e)=e'$.
\end{enumerate}
% The definition ends at the bottom of the manuscript page.
.
% Faithful translation retaining the manuscript's sequence and notation.
% No preamble; start a new lecture on a new page.

\clearpage
\subsection*{Lecture 12 - Page 1}

\begin{center}
\fbox{\textbf{LECTURE 12}}
\end{center}
\begin{flushright}
\textit{3-Jan-1995}
\end{flushright}

\medskip
\noindent\underline{Generalities on monoids}

\bigskip
\noindent\fbox{T.}\quad Let $(M,\cdot,e)$ be a monoid, $\varnothing\neq S\subseteq M$,
\[
\langle S\rangle
=\{x_1\cdots x_n\mid x_i\in S,\ n\geq 1\}\cup\{e\}.
\]
\noindent
\begin{enumerate}
\item[(1)] $\langle S\rangle$ is a submonoid of $M$, $\langle S\rangle\leq M$.
\item[(2)] $N\leq M$, $S\subseteq N\ \Longrightarrow\ \langle S\rangle\subseteq N$.
\end{enumerate}

\noindent A non-empty subset $S\subseteq M$ is called a \underline{submonoid}
if $S$ is closed under the operation and $e\in S$.

\bigskip
\noindent\underline{Proof.} Evident.

\bigskip
\noindent\underline{Remark.}\quad
$\langle S\rangle$ is the smallest submonoid of $M$ containing $S$.

\medskip
\noindent $\langle S\rangle$ is called the submonoid of $M$ generated by $S$.

\bigskip
\noindent\underline{Definition.}\quad
Let $M$ be a monoid. A subset $S\subseteq M$ is called a
\underline{generating set} for $M$ if
\[
\langle S\rangle=M.
\]

\bigskip
\noindent\underline{Example.}\quad
For $(\mathbb{N}^{*},\cdot,1)$, a generating set is the set of prime numbers,
or the set of prime numbers together with $\{1\}$.
% The handwritten example reads “{prime numbers} or {prime numbers}+{1}”.
% The final marks at the line's end are unclear; no additional example supplied.

\bigskip
\noindent\underline{Definition.}\quad
A monoid $M$ is called \underline{finitely generated}
if its generating set is finite.

\medskip
\noindent\underline{Remark.}\quad $M=\langle M\rangle$.

\bigskip
\noindent $R$ - a congruence on $M$
(an equivalence relation compatible with $\cdot$).
\[
\frac{M}{R}=\{\hat a\mid a\in M\},
\qquad \hat a\,\hat b=\widehat{ab},
\qquad \left(\frac{M}{R},\cdot,\hat e\right)
\text{ - the quotient monoid.}
\]

\bigskip
\noindent\underline{Example.}\quad
$(\mathbb{Z},\cdot,1)$, $R_n$ - congruence modulo $n$,
\[
\frac{\mathbb{Z}}{R_n}=(\mathbb{Z}_n,\cdot,\hat 1).
\]
% The manuscript uses equality here to identify the quotient with the quotient monoid.

\bigskip
\noindent\underline{Definition.}\quad
Let $(M,\cdot,e)$ and $(M',\cdot,e')$ be two monoids.
A map $f:M\to M'$ is called a \underline{monoid homomorphism} if:
\begin{enumerate}
\item[(1)] $f(xy)=f(x)f(y)$, $(\forall)\,x,y\in M$.
\item[(2)] $f(e)=e'$.
\end{enumerate}
% The definition ends at the bottom of the manuscript page.
.
% Faithful translation retaining the manuscript's sequence and notation.
% No preamble; start a new lecture on a new page.

\clearpage
\subsection*{Lecture 12 - Page 1}

\begin{center}
\fbox{\textbf{LECTURE 12}}
\end{center}
\begin{flushright}
\textit{3-Jan-1995}
\end{flushright}

\medskip
\noindent\underline{Generalities on monoids}

\bigskip
\noindent\fbox{T.}\quad Let $(M,\cdot,e)$ be a monoid, $\varnothing\neq S\subseteq M$,
\[
\langle S\rangle
=\{x_1\cdots x_n\mid x_i\in S,\ n\geq 1\}\cup\{e\}.
\]
\noindent
\begin{enumerate}
\item[(1)] $\langle S\rangle$ is a submonoid of $M$, $\langle S\rangle\leq M$.
\item[(2)] $N\leq M$, $S\subseteq N\ \Longrightarrow\ \langle S\rangle\subseteq N$.
\end{enumerate}

\noindent A non-empty subset $S\subseteq M$ is called a \underline{submonoid}
if $S$ is closed under the operation and $e\in S$.

\bigskip
\noindent\underline{Proof.} Evident.

\bigskip
\noindent\underline{Remark.}\quad
$\langle S\rangle$ is the smallest submonoid of $M$ containing $S$.

\medskip
\noindent $\langle S\rangle$ is called the submonoid of $M$ generated by $S$.

\bigskip
\noindent\underline{Definition.}\quad
Let $M$ be a monoid. A subset $S\subseteq M$ is called a
\underline{generating set} for $M$ if
\[
\langle S\rangle=M.
\]

\bigskip
\noindent\underline{Example.}\quad
For $(\mathbb{N}^{*},\cdot,1)$, a generating set is the set of prime numbers,
or the set of prime numbers together with $\{1\}$.
% The handwritten example reads “{prime numbers} or {prime numbers}+{1}”.
% The final marks at the line's end are unclear; no additional example supplied.

\bigskip
\noindent\underline{Definition.}\quad
A monoid $M$ is called \underline{finitely generated}
if its generating set is finite.

\medskip
\noindent\underline{Remark.}\quad $M=\langle M\rangle$.

\bigskip
\noindent $R$ - a congruence on $M$
(an equivalence relation compatible with $\cdot$).
\[
\frac{M}{R}=\{\hat a\mid a\in M\},
\qquad \hat a\,\hat b=\widehat{ab},
\qquad \left(\frac{M}{R},\cdot,\hat e\right)
\text{ - the quotient monoid.}
\]

\bigskip
\noindent\underline{Example.}\quad
$(\mathbb{Z},\cdot,1)$, $R_n$ - congruence modulo $n$,
\[
\frac{\mathbb{Z}}{R_n}=(\mathbb{Z}_n,\cdot,\hat 1).
\]
% The manuscript uses equality here to identify the quotient with the quotient monoid.

\bigskip
\noindent\underline{Definition.}\quad
Let $(M,\cdot,e)$ and $(M',\cdot,e')$ be two monoids.
A map $f:M\to M'$ is called a \underline{monoid homomorphism} if:
\begin{enumerate}
\item[(1)] $f(xy)=f(x)f(y)$, $(\forall)\,x,y\in M$.
\item[(2)] $f(e)=e'$.
\end{enumerate}
% The definition ends at the bottom of the manuscript page.
